% BHU PYQ -- Physics Mechanics & General Properties of Matter (End Term) (2024-25 NEP)
\documentclass[11pt,a4paper]{article}
\usepackage[a4paper,top=2.5cm,bottom=2.5cm,left=2.8cm,right=2.8cm,headheight=14pt]{geometry}
\usepackage{amsmath,amssymb}
\usepackage[T1]{fontenc}
\usepackage[utf8]{inputenc}
\usepackage{lmodern,microtype,fancyhdr,enumitem,hyperref,lastpage,parskip}

\hypersetup{colorlinks=true,urlcolor=black,linkcolor=black,pdftitle=PHYMJ11: Mechanics & General Properties of Matter (End Term) -- BHU NEP 2024-25}

\pagestyle{fancy}\fancyhf{}
\renewcommand{\headrulewidth}{0.4pt}\renewcommand{\footrulewidth}{0.4pt}
\fancyhead[L]{\small\textit{Banaras Hindu University}}
\fancyhead[R]{\small\textit{PHYMJ11: Mechanics & General Properties of Matter (End Term) (2024-25)}}
\fancyfoot[C]{\small Page \thepage\ of \pageref{LastPage}}

\newlist{parts}{enumerate}{2}
\setlist[parts,1]{label=(\alph*),leftmargin=*,itemsep=4pt,topsep=3pt}
\setlist[parts,2]{label=(\roman*),leftmargin=*,itemsep=2pt,topsep=2pt}
\newcommand{\pts}[1]{\hfill\textit{[#1]}}

\begin{document}
\begin{center}
  {\large\bfseries BANARAS HINDU UNIVERSITY}\\[2pt]
  {\normalsize U.G. Semester 1 (NEP) Examination, 2024-25}\\[6pt]
  \rule{0.8\linewidth}{0.5pt}\\[6pt]
  {\large\bfseries PHYSICS}\\[4pt]
  {\large\bfseries Paper: PHYMJ11 --- Mechanics & General Properties of Matter (End Term)}\\[4pt]
  {\normalsize Time: 2:30 Hours \hfill Full Marks: 70}
\end{center}
\rule{\linewidth}{0.4pt}
\smallskip
\noindent\textit{\noindent\textit{Instructions: Attempt \textbf{FOUR} questions including Question No.~1, which is compulsory. The figures in the right-hand margin indicate marks.}
\bigskip

\noindent\textbf{Question 1.} \textit{(Compulsory --- Attempt any four)}
\begin{parts}
  \item Show that in a perfectly inelastic collision in a laboratory system, there is always a loss of kinetic energy.\pts{3.5}
  \item What couple must be applied to a one-meter-long wire, having $1\text{ mm}$ diameter, in order to twist one end through $90^{\circ}$, while the other end is fixed? Rigidity of the material is $\eta = 2.8 \times 10^{10}\text{ N/m}^2$.\pts{3.5}
  \item A uniform circular ring of mass $M$ and radius $R$ hangs in a vertical plane supported by a knife edge at one point on its inner circumference. Calculate the natural frequency of small oscillations.\pts{3.5}
  \item Explain inertial and non-inertial frames of reference with suitable examples.\pts{3.5}
  \item If the centre of mass of three particles of masses $1\text{ kg}$, $2\text{ kg}$, and $3\text{ kg}$ is at the point $(3, 3, 3)$, where should a fourth mass of $4\text{ kg}$ be placed so that the centre of mass of the four particles is at $(1, 1, 1)$?\pts{3.5}
  \item Given that the orbital velocity of the Sun about the centre of our galaxy is $3 \times 10^7\text{ cm/s}$ and its distance is nearly $3 \times 10^{22}\text{ cm}$ from the galactic axis, estimate the mass of the galaxy ($G = 6.67 \times 10^{-8}\text{ CGS units}$).\pts{3.5}
\end{parts}

\bigskip
\noindent\textbf{Question 2.}
\begin{parts}
  \item Derive an expression for the twisting couple per unit twist for a solid cylinder. Show that for cylinders of the same mass, material, and length, the twisting couple per unit twist is greater for a hollow cylinder than for a solid one.\pts{7}
  \item Derive the relation connecting Young's modulus ($Y$), Bulk modulus ($K$), and Shear modulus ($\eta$) of an isotropic elastic material.\pts{7}
\end{parts}

\bigskip
\noindent\textbf{Question 3.}
\begin{parts}
  \item Write down the differential equation of a damped harmonic oscillator and solve it to explain three different conditions of damping (overdamped, critically damped, underdamped).\pts{7}
  \item The frequency of an undamped harmonic oscillator is equal to twice the frequency experienced by the oscillator with damping. Calculate the logarithmic decrement of the system.\pts{7}
\end{parts}

\bigskip
\noindent\textbf{Question 4.}
\begin{parts}
  \item State and prove Bernoulli's theorem for fluid flow: $\frac{1}{2}v^2 + \frac{P}{\rho} + gh = \text{constant}$.\pts{14}
\end{parts}


\vfill\rule{\linewidth}{0.4pt}
\end{document}
